\section[拓扑场论、界面与边界条件]{拓扑场论、界面与边界条件（4 月 8 日）}

\srcpage{57}

TODO：我迟到了几分钟（四分钟）。

我们要从拓扑场论的角度谈算子与缺陷。先定义（定向）配边 $2$-范畴 $\mathrm{Bord}_2^{\mathrm{or}}$：对象是带二维领邻域的闭定向 $0$-流形；$1$-态射是配边（紧致、定向、带一维领邻域）；$2$-态射是带角的紧定向 $2$-流形，它们是这些 $1$-态射之间的配边。对称幺半结构由不交并给出。

配边成族出现，或者等价地说有分类空间。例如若 $\Sigma_g$ 是亏格 $g$ 的闭连通定向曲面，可以把它看成从空 $1$-流形（本身是从空 $0$-流形到自身的配边）到自身的配边；则与 $\Sigma_g$ 微分同胚的曲面族由到分类空间 $B\mathrm{Diff}(\Sigma_g)$ 的映射的同伦类分类，而该分类空间同胚于亏格 $g$ 曲线的模空间 $\cM_g$。

二维拓扑场论于是是对称幺半函子

\[Z:\mathrm{Bord}_2^{\mathrm{or}}\longrightarrow(\cC,\otimes,1),\]

取值到某个对称幺半 $2$-范畴。我们希望目标 $2$-范畴是 $\C$-线性的（大致说，hom-范畴是 $\C$-线性范畴），并希望它扩张通常的 TFT 概念，使得 $\Omega\cC\simeq\Vect$。这里 $\Omega(-):=\Hom_{(-)}(1,1)$。$\mathrm{Bord}_2^{\mathrm{or}}$ 的幺半单位是空 $0$-流形，所以 $\Omega\mathrm{Bord}^{\mathrm{or}}$ 是通常（非扩张）的定向 $2$-流形配边范畴，而扩张 TFT 限制为通常 TFT。

目标 $\cC$ 的一个「万有」选择是 $\C$-线性范畴的 $2$-范畴 $\mathrm{Cat}_\C$：若 $\cC$ 是另一个满足我们要求的范畴，则（TODO：怎么做到？）可以从 $Z(pt_+)$ 提取一个 $\C$-线性范畴，于是映射 $\mathrm{Bord}_2^{\mathrm{or}}\to\cC$ 延拓为 $\cC\to\mathrm{Cat}_\C$。实践中 $\cC$ 常等价或几乎等价于 $\mathrm{Cat}_\C$。

\begin{example}
另一个可作目标的是 Morita $2$-范畴 $\mathrm{Alg}_\C$：对象是 $\C$ 上结合代数；态射 $A\to B$ 是 $(B,A)$-双模；$2$-态射是双模同态。
\end{example}

有对称幺半函子 $\mathrm{Alg}_\C\to\mathrm{Cat}_\C$，把代数送到其表示范畴，并把 $\mathrm{Alg}_\C$ 嵌入为 $\mathrm{Cat}_\C$ 的满子范畴，所以这两个选择其实差别不大。

我们用 TFT 语言解释这套 $2$-范畴的东西。上次定义了 TFT 之间的界面；人们定义二部 $2$-流形的配边范畴（曲面沿余维一轨迹分解为两个不同子流形，染红与染蓝）。给定两个 TFT $Z,W$，可以定义从部分二部配边范畴到 $\mathrm{Cat}_\C$ 的函子，把纯红曲面送到 $Z$、纯蓝曲面送到 $W$；但在同时含两部分的曲面上还需要更多数据才能定义。一个称为\srcpage{58}带奇点的协边假设（cobordism hypothesis with singularities）的定理——由 Lurie 证明——说明这个额外数据恰是一个界面，即 $\cC$ 中的映射 $Z(pt)\to W(pt)$。\footnote{这里至少在概念上，重要的是把二部范畴中所有对象与态射连同其领邻域一起画出来。因此，例如有一个点，其领邻域左侧染红、右侧染蓝，下面映射 $Z(pt_+)\to W(pt_+)$ 就是这样得到的。}

这是引出下述想法的冗长动机：若令 $W$ 为平凡理论（取值于单位的常函子），则从 $W$ 到 $Z$ 的界面等同于 $\Hom(1,Z(pt_+))$ 的对象，也等同于 $Z(pt_+)$ 的对象。用 TFT 语言，这定义了 $Z$ 的「边界理论」：可以把 $Z$ 扩张到带标记边界分支的 $2$-流形配边范畴。物理上，可以想把生活在「体理论」边界上的场论想成某种耦合系统。而这些正是 $Z(pt_+)$ 的对象。

这套东西有点令人困惑——$2$d TFT 的边界条件自然会构成范畴，即使没有固定 $\cC=\mathrm{Cat}_\C$，这可以看成「底层范畴」操作 $\cC\to\mathrm{Cat}_\C$。但有时这个映射不太有趣或是平凡的。不过实践中，对我们关心的 $\cC$ 的选择，结果都不错。

也可以推广到余维 $k$ 缺陷理论：把 $Z$ 推广到带余维 $k$ 嵌入标记子流形的 $n$-流形。$k=1$ 时这与边界条件或界面并不完全相同，因为并非所有余维一子流形都分离；我们有更多灵活性，因为可以在分离子流形两侧改变 TFT。余维一缺陷由 $Z$ 的自同态数据给出（从 $Z$ 到自身的界面是映射 $Z\to Z$）。这里的差别在于 $S^0$ 不连通！

\begin{remark}
同伦论中的 Bass–Sullivan 理论包含一个定理：所有广义上同调理论都可实现为带某类奇点的流形的配边理论。（有些不需要奇点，例如球谱表示稳定上同伦，由标架配边得到。）带奇点的协边假设是这件事的范畴化版本：我们允许某类「奇点」（缺陷），并得到 $(\infty,n)$-范畴中所有（足够小，即完全可对偶的）对象、它们之间的态射（来自界面）等等。
\end{remark}

我们看看 $2$d TFT $Z$ 的边界条件如何构成范畴。由抽象废话我们知道这是对的，但我们要几何地看：界面之间还可以有界面！例如可以为 $Z$ 选择两个不同的边界条件 $B$ 与 $R$；从 $B$ 到 $R$ 的界面允许我们在如图 2 的图景上定义 TFT。

原文图 2 画的是「在边界条件 $B$ 与 $R$ 之间给定界面 $\phi$ 后可求值的配边」；因图片未复刻，正文已叙述其含义。\yizhu{原文图 2 的图注为「A bordism that can be evaluated given an interface $\phi$ between the boundary conditions $B$ and $R$」；此处从略。}{A（原文 p.58）}

那个图不闭合，但我们可以把它封上，得到带角的两条带标记边界的 $1$-流形之间的配边，它们被赋予向量空间。所以从 $B$ 到 $R$ 的界面集合构成向量空间 $Z([0,1])$，我们定义 $\Hom(B,R)$ 为这个向量空间。

这些界面可以复合，靠的是「双帽」配边（TODO）定义线性映射 $\Hom(B,R)\otimes\Hom(R,G)\to\Hom(B,G)$，其中 $G$ 是第三个界面。

同样的想法在任意维 $n$ 都行：可以构造边界条件的 $(n-1)$-范畴。高阶范畴论正是在这里物理地出现：人们研究缺陷之间的缺陷之间的缺陷。

\begin{example}[有限规范场论]
这个理论有许多名字，包括（未扭）Dijkgraaf–Witten 理论与有限群 Yang–Mills 理论。

取有限群 $G$；我们通过线性化流形上 $G$-局部系统的空间来定义 $2$d TFT $Z_G:\mathrm{Bord}_2^{\mathrm{or}}\to\mathrm{Cat}_\C$。某种意义上这是 Lagrange 场论：给流形附上场的空间（嗯，群胚），再通过取层线性化，得到范畴。这正是我们\srcpage{59}在余维 $2$ 附上范畴的动机：我们希望得到场空间上的层，而不只是某个抽象函子；而这里确实得到场空间（群胚）上的层。

从闭流形 $\Sigma$ 开始。则 $Z_G(\Sigma)$ 是 $\mathrm{Loc}_G(\Sigma)$ 的基数。正确的计数方式是按自同构加权：$\mathrm{Loc}_G(\Sigma)$ 是自同构群有限的群胚，所以可定义和

\[Z_G(\Sigma)=\sum_{[P]\in\pi_0(\mathrm{Loc}_G(\Sigma))}\frac{1}{\#\Aut(P)}.\]

我们不太关心这些数，而更关心更低余维上发生的事。对闭 $1$-流形 $N$，附上 $\mathrm{Loc}_G(N)$ 上的函数向量空间：

\[Z_G(N):=\C[\mathrm{Loc}_G(N)].\]

例如 $S^1$ 上的局部系统等同于群胚 $G/G$（$G$ 通过共轭作用在自身上），而 $G/G$ 上的函数等同于类函数 $\C[G]^G$。

对 $0$-流形 $P$，附上 $\mathrm{Loc}_G(P)$ 上的层范畴。

配边上如何？从 $N_1$ 到 $N_2$ 的配边 $M$ 给出图表

\begin{tikzcd}
 & \mathrm{Loc}_G(M) \arrow[dl] \arrow[dr] & \\
 \mathrm{Loc}_G(N_1) & & \mathrm{Loc}_G(N_2)
\end{tikzcd}

由于 $G$ 有限，可以对函数或层做推拉操作，得到配边上的映射。

从这个观点，$G$ 的表示是什么？$Z_G(pt)=\Vect(pt/G)=\Rep_G$，所以表示就是边界条件。确切地说，给定 $G$-表示 $V$，我们得到把理论扩张到带标记红色边界分支的流形上的数据。例如区间 $[0,1]$ 带一个红色边界分支，取值于对象 $V\in Z_G(pt)$；而圆柱 $[0,1]\times S^1$ 定义 $Z_G(S^1)$ 的一个元素，结果恰是 $V$ 的特征标。\footnote{若你熟悉 TFT 中「与 $S^1$ 取积是一种迹」的想法，这就很合理。}

物理上，把 $V\in\Rep_G$ 想成某个（拓扑）量子力学的态空间。$G$-作用是额外数据，描述如何把它耦合到体 $2$d 理论 $Z_G$。例如正则表示（$V=\C[G]$）称为 $Z_G$ 的 Dirichlet 边界条件，而 $V=\C$ 称为 $Z_G$ 的 Neumann 边界条件。

给定两个边界条件（$G$-表示）$V,W$，它们之间能放什么界面？我们知道这是 $\Hom_{\Rep_G}(V,W)$。复合由缺陷碰撞实现（把它们放近，然后看成一个缺陷）。

也可以用局部系统来想：若 $G$ 作用在空间 $X$ 上，它也作用在 $\C[X]$ 上。带红色边界标号 $X$ 的 $2$-流形 $\Sigma$ 上的场空间，是 $G$-局部系统连同边界上相伴 $X$-丛的一个截面。换种说法，体内有到群胚 $pt/G$ 的映射，边界上把这个映射沿 $X/G\to pt/G$ 提升。于是 Dirichlet 边界条件 $X=G$ 给出的场空间是边界上带平凡化的局部系统。
\end{example}

下次我们会进一步讨论这个例子，例如子群 $K\subset G$ 给出的 $Z_G$ 的边界条件 $G/K$。双陪集空间 $K\backslash G/K$ 在这个框架中有自然解释，通过从 $\C[G/K]$ 到自身的界面：相应配边上的场是 $G$-局部系统，并提升为边界上的 $K$-局部系统，这意味着映射到 $K\backslash G/K$。
